%!TEX root = ../main.tex
% Declaración sobre el uso de inteligencia artificial.
% Exigencia de la cátedra (clase 5): el uso de IA está permitido y
% fomentado, pero debe declararse. Esta página resume el uso; el
% registro de prompts va en el apéndice~\ref{apx:prompts}.

\begin{guia}[Guía de la cátedra: política de uso de IA]
La IA es parte del estado del arte de la industria y su uso está
permitido, con condiciones:
\begin{enumerate}
    \item \textbf{Declararlo}: acá y en la Metodología (qué tareas fueron
          asistidas y con qué herramientas).
    \item \textbf{Autoría}: las ideas, decisiones y conclusiones son del
          estudiante. La IA asiste; no decide.
    \item \textbf{Verificación}: toda cifra, cita o referencia sugerida por
          una IA se confirma en la fuente original antes de usarla.
    \item \textbf{Trazabilidad}: los prompts relevantes se guardan en una
          carpeta compartida con la cátedra.
    \item \textbf{Comprensión}: no se incorpora texto que no se haya leído
          y entendido. Plagio, incluso si viene de un modelo de lenguaje,
          es plagio.
\end{enumerate}
\end{guia}

En este trabajo se utilizaron herramientas de inteligencia artificial
generativa tanto en el desarrollo del sistema como en la elaboración del
documento. La tabla~\ref{tab:herramientas-ia} resume cada herramienta y
las tareas que asistió; el método de desarrollo se describe en el
capítulo~\ref{cap:metodologia} (sección~\ref{sec:desarrollo-ia}).

\begin{table}[H]
    \centering
    \caption{Herramientas de IA generativa utilizadas}
    \label{tab:herramientas-ia}
    \small
    \begin{tabularx}{\textwidth}{>{\raggedright\arraybackslash}p{4.2cm} L}
        \toprule
        \tb{Herramienta} & \tb{Tareas asistidas} \\
        \midrule
        OpenCode (modelo Big Pickle) & Implementación del código del \gls{mvp} a partir de la Constitución, las especificaciones, los planes técnicos y la lista de tareas redactados previamente, en las primeras etapas del desarrollo \\
        Gemini (Gemini 3.1 Pro) & Consultas esporádicas, como segunda opinión sobre decisiones puntuales \\
        ChatGPT (modelo predeterminado) & Consultas esporádicas, como tercera mirada sobre decisiones puntuales \\
        Claude Code (Claude Opus 5, Claude Opus 5.5 y Claude Sonnet 5) & Revisión de la consistencia del documento, relevamiento de tarifas y fuentes, recálculo y edición del modelo financiero, y redacción asistida de capítulos y apéndices \\
        \bottomrule
    \end{tabularx}
\end{table}

El uso se rigió por los siguientes criterios:

\begin{itemize}
    \item \textbf{Autoría.} La definición del problema, la arquitectura,
          el alcance, los criterios de evaluación y las decisiones del
          proyecto son del autor. Las herramientas propusieron
          alternativas y redacciones; la elección entre ellas fue siempre
          del autor, y las decisiones de arquitectura se documentaron como
          registros de decisión validados por él.
    \item \textbf{Verificación.} Toda tarifa, cotización y dato de
          mercado se confirmó en la fuente original, que se cita en cada
          caso. Los indicadores financieros se calcularon en la planilla
          del anexo~\ref{anx:modelo-financiero} y se replicaron de forma
          independiente antes de incluirlos en el texto. El código se
          verificó con las pruebas automatizadas del repositorio y la
          integración continua.
    \item \textbf{Trazabilidad.} Las instrucciones con las que se
          desarrolló el sistema están versionadas en el propio
          repositorio (anexo~\ref{anx:repositorios}). Los prompts más
          relevantes de la elaboración del documento se registran en el
          apéndice~\ref{apx:prompts}.
    \item \textbf{Comprensión.} El autor leyó, revisó y adaptó todo el
          contenido asistido, y asume la responsabilidad total sobre el
          contenido de este documento.
\end{itemize}
